\documentclass[exercice]{thedocument}%
\startdatakeys
\source{}
\BO{2026}
\cycle{4}
\competencesDuSocle{CA,RA}
\objectifsApprentissage{B2eb,B2ec}
\automatismes{B3cb}
\enddatakeys
\begin{document}%
Grâce à la lecture de photos satellites, des chercheurs brésiliens ont pu
calculer les longueurs de l'Amazone et du Nil avec une très grande précision.
Ils ont trouvé 6\,992,06 km pour l'Amazone et 6\,852,15 km pour le Nil. Ainsi,
l'Amazone est le plus long fleuve du monde.\vskip10pt De combien de km
dépasse-t-il le Nil ?
\end{document}
